% Figure: the Carnot cycle on a T-S diagram, where it is a rectangle.
\begin{figure}[htbp]
\centering
\begin{tikzpicture}
\begin{axis}[
  width=10cm, height=7.5cm,
  xlabel={entropy $S$ (arb.)}, ylabel={temperature $T$ (arb.)},
  xmin=0.5, xmax=4.5, ymin=0, ymax=6,
  axis lines=left,
  tick label style={font=\scriptsize},
  label style={font=\small},
  clip=false,
]
% Carnot rectangle: T_h=5, T_c=2, S from 1 to 4.
% 1->2 isothermal expansion at T_h: S increases 1->4
\addplot[engred, very thick] coordinates {(1,5) (4,5)};
% 2->3 adiabatic (isentropic) expansion: S=4 constant, T 5->2
\addplot[englime!50!black, very thick] coordinates {(4,5) (4,2)};
% 3->4 isothermal compression at T_c: S decreases 4->1
\addplot[engblue, very thick] coordinates {(4,2) (1,2)};
% 4->1 adiabatic compression: S=1 constant, T 2->5
\addplot[engorange, very thick] coordinates {(1,2) (1,5)};
% shade the enclosed area
\addplot[fill=engblue!10, draw=none] coordinates {(1,2) (4,2) (4,5) (1,5)} \closedcycle;
% Qh = area under top edge down to T=0
\node[engred, font=\scriptsize] at (axis cs:2.5,5.3) {$Q_h = T_h \Delta S$};
\node[engblue, font=\scriptsize] at (axis cs:2.5,1.6) {$Q_c = T_c \Delta S$};
\node[font=\scriptsize] at (axis cs:2.5,3.5) {net work $= (T_h-T_c)\Delta S$};
% state labels
\filldraw[engblack] (axis cs:1,5) circle [radius=2pt];
\node[font=\scriptsize, anchor=south east] at (axis cs:1,5) {1};
\filldraw[engblack] (axis cs:4,5) circle [radius=2pt];
\node[font=\scriptsize, anchor=south west] at (axis cs:4,5) {2};
\filldraw[engblack] (axis cs:4,2) circle [radius=2pt];
\node[font=\scriptsize, anchor=north west] at (axis cs:4,2) {3};
\filldraw[engblack] (axis cs:1,2) circle [radius=2pt];
\node[font=\scriptsize, anchor=north east] at (axis cs:1,2) {4};
\end{axis}
\end{tikzpicture}
\caption[The Carnot cycle on a $T$-$S$ diagram]{The same Carnot cycle drawn on
temperature-entropy axes is a rectangle. Heat exchanged along any reversible
path is the area under the path, $Q = \int T\,dS$, so the heat absorbed at the
top edge is $Q_h = T_h\,\Delta S$ and the heat rejected at the bottom edge is
$Q_c = T_c\,\Delta S$. The two adiabatic legs are vertical because a reversible
adiabatic process holds entropy constant. The enclosed rectangle is the net
work $(T_h - T_c)\,\Delta S$, and the efficiency $1 - T_c/T_h$ reads directly
off the ratio of the two horizontal edges' heights.}
\label{fig:vol04ch01:carnot-ts}
\end{figure}
